\documentclass[12pt, a4paper]{ctexart}
\usepackage[margin=2cm]{geometry}
\usepackage{amsmath, amssymb}
\usepackage{enumitem}
\usepackage{fancyhdr}
\usepackage{xcolor}
\usepackage{tikz}
\usepackage{graphicx}
\usepackage{tabularx}
\usepackage{multicol}

\setCJKmainfont{Hiragino Sans GB}
\setCJKsansfont{Heiti SC}

\pagestyle{fancy}
\fancyhf{}
\fancyhead[L]{培优专题·压轴题模型}
\fancyhead[R]{测试卷}
\fancyfoot[C]{\thepage}
\renewcommand{\headrulewidth}{0.4pt}

\begin{document}

\begin{center}
{\Large\bfseries 培优专题·压轴题模型 测试卷}\\[0.5em]
{\normalsize 满分150分\quad 考试时间120分钟}
\end{center}

\vspace{0.5em}

\noindent\textbf{一、单项选择题（本大题共8小题，每小题5分，共40分。每小题只有一个选项符合题意。）}

\begin{enumerate}[label=\Alph*.]

\item 已知函数$f(x)=e^x-x-1$，则$f(x)$在$\mathbb{R}$上（\quad）
\begin{enumerate}[label=]
\item 有唯一零点$x=0$ \qquad \item 有两个零点
\item 无零点 \qquad \item 有三个零点
\end{enumerate}

\item 设$a=\ln 2$，$b=\dfrac{1}{e}$，$c=2^{-\frac{1}{2}}$，则$a$、$b$、$c$的大小关系为（\quad）
\begin{enumerate}[label=]
\item $a>b>c$ \qquad \item $b>a>c$ \qquad \item $c>a>b$ \qquad \item $a>c>b$
\end{enumerate}

\item 已知函数$f(x)=x^3-3x^2+ax$在$\mathbb{R}$上单调递增，则实数$a$的取值范围是（\quad）
\begin{enumerate}[label=]
\item $a\geq 3$ \qquad \item $a>3$ \qquad \item $a\geq 0$ \qquad \item $a>0$
\end{enumerate}

\item 若不等式$x\ln x\geq x-1$对一切$x>0$恒成立，则下列结论正确的是（\quad）
\begin{enumerate}[label=]
\item $\ln 2>\dfrac{2}{3}$ \qquad \item $\ln 3>\dfrac{3}{4}$ \qquad \item $\ln 5>\dfrac{4}{5}$ \qquad \item $\ln e>1$
\end{enumerate}

\item 函数$f(x)=\dfrac{\ln x}{x}$（$x>0$）的极大值为（\quad）
\begin{enumerate}[label=]
\item $1$ \qquad \item $\dfrac{1}{e}$ \qquad \item $e$ \qquad \item $\dfrac{e}{2}$
\end{enumerate}

\item 已知数列$\{a_n\}$满足$a_1=1$，$a_{n+1}=a_n+\dfrac{1}{a_n}$，则$a_{10}$的整数部分为（\quad）
\begin{enumerate}[label=]
\item $3$ \qquad \item $4$ \qquad \item $5$ \qquad \item $6$
\end{enumerate}

\item 已知函数$f(x)=x-e^x+1$，则不等式$f(\ln x)<0$的解集为（\quad）
\begin{enumerate}[label=]
\item $(0,1)$ \qquad \item $(1,+\infty)$ \qquad \item $(0,e)$ \qquad \item $(e,+\infty)$
\end{enumerate}

\item 设函数$f(x)=\ln x-ax^2+bx$（$a>0$），若$f(x)$在$x=1$处取得极值，则$\dfrac{b}{a}$的值为（\quad）
\begin{enumerate}[label=]
\item $1$ \qquad \item $2$ \qquad \item $3$ \qquad \item $4$
\end{enumerate}

\end{enumerate}

\vspace{0.5em}

\noindent\textbf{二、多项选择题（本大题共4小题，每小题5分，共20分。全部选对得5分，部分选对得2分，有选错的得0分。）}

\begin{enumerate}[label=\arabic*.]

\item 已知函数$f(x)=x^3-12x$，下列结论正确的是（\quad）
\begin{enumerate}[label=]
\item $f(x)$在$(-2,2)$上单调递减
\item $f(x)$的极大值为$16$
\item $f(x)$的极小值为$-16$
\item 方程$f(x)=0$有三个不同实数根
\end{enumerate}

\item 关于函数$g(x)=e^x-x-1$，下列结论正确的是（\quad）
\begin{enumerate}[label=]
\item $g(x)\geq 0$对一切$x\in\mathbb{R}$恒成立
\item $g(x)$在$x=0$处取得最小值$0$
\item 当$x>0$时，$e^x>x+1$
\item 当$x<0$时，$e^x<x+1$
\end{enumerate}

\item 已知$a>0$，$b>0$，$a+b=1$，下列不等式正确的是（\quad）
\begin{enumerate}[label=]
\item $ab\leq\dfrac{1}{4}$
\item $\dfrac{1}{a}+\dfrac{1}{b}\geq 4$
\item $\sqrt{a}+\sqrt{b}\leq\sqrt{2}$
\item $a^2+b^2\geq\dfrac{1}{2}$
\end{enumerate}

\item 已知函数$f(x)=\ln(1+x)-x+\dfrac{x^2}{2}$，下列结论正确的是（\quad）
\begin{enumerate}[label=]
\item $f(x)$在$(-1,0)$上单调递减
\item $f(x)$在$(0,+\infty)$上单调递增
\item $f(x)\geq 0$对$x>-1$恒成立
\item $\ln(1+x)\leq x-\dfrac{x^2}{2}$（$x>-1$）
\end{enumerate}

\end{enumerate}

\vspace{0.5em}

\noindent\textbf{三、填空题（本大题共4小题，每小题5分，共20分。）}

\begin{enumerate}[label=\arabic*.]

\item 函数$f(x)=x^3-3x+1$在区间$[-3,2]$上的最大值与最小值之差为 \underline{\hspace{3cm}}。

\item 已知$f(x)=\ln x+x$，$g(x)=e^x+x$，若$f(a)=g(b)=2$，则$a+b=$ \underline{\hspace{3cm}}。

\item 已知数列$\{a_n\}$满足$a_1=1$，$a_{n+1}=\dfrac{a_n}{1+a_n}$，则$a_n=$ \underline{\hspace{3cm}}。

\item 已知函数$f(x)=e^x-ax^2$（$a>0$），若$f(x)\geq 0$对$x\geq 0$恒成立，则$a$的最大值为 \underline{\hspace{3cm}}。

\end{enumerate}

\vspace{0.5em}

\noindent\textbf{四、解答题（本大题共6小题，共70分。解答应写出文字说明、证明过程或演算步骤。）}

\begin{enumerate}[label=\arabic*.]

\item （本小题满分10分）\\
已知函数$f(x)=x\ln x$。\\
（1）求$f(x)$的单调区间和极值；\\
（2）证明：$\ln x>\dfrac{2(x-1)}{x+1}$（$x>1$）。
\vspace{6cm}

\item （本小题满分12分）\\
已知函数$f(x)=e^x-ax$（$a>0$）。\\
（1）求$f(x)$的最小值$g(a)$；\\
（2）求$g(a)$的最大值。
\vspace{6cm}

\item （本小题满分12分）\\
已知数列$\{a_n\}$满足$a_1=2$，$a_{n+1}=a_n^2-na_n+1$。\\
（1）求$a_2$、$a_3$、$a_4$的值；\\
（2）猜想$a_n$关于$n$的表达式，并用数学归纳法证明；\\
（3）求$a_1+a_2+\cdots+a_{10}$。
\vspace{6cm}

\item （本小题满分12分）\\
已知函数$f(x)=\ln x-ax+\dfrac{a}{x}-1$（$a\in\mathbb{R}$）。\\
（1）当$a=1$时，求$f(x)$的单调区间；\\
（2）当$a>0$时，讨论$f(x)$的零点个数；\\
（3）若$f(x)\leq 0$对$x>0$恒成立，求$a$的取值范围。
\vspace{6cm}

\item （本小题满分12分）\\
已知函数$f(x)=\ln x$，$g(x)=\dfrac{a}{x}$（$a>0$）。\\
（1）求$f(x)$与$g(x)$的图象的公切线方程；\\
（2）证明：$\ln x\leq\dfrac{x^2-1}{2x}$（$x>0$）；\\
（3）证明：$\dfrac{1}{\ln 2}+\dfrac{1}{\ln 3}+\cdots+\dfrac{1}{\ln 100}<\dfrac{99}{2}$。
\vspace{6cm}

\item （本小题满分12分）\\
已知函数$f(x)=e^x-\dfrac{1}{2}x^2-x$。\\
（1）求$f(x)$的单调区间；\\
（2）证明：$\displaystyle\sum_{k=1}^{n}\dfrac{1}{k!}<\dfrac{e}{e-1}$；\\
（3）证明：$e^{\pi}>\pi^e$。
\vspace{6cm}

\end{enumerate}

\end{document}
